\begin{frame}{왜 이렇게 느리게 수렴하는가}
  \begin{block}{힌트}
    $p_0$ 가 1에 가까워질수록:
    \begin{itemize}
      \item 그래디언트 성분 $(1-p_0)$ 이 \textbf{줄어}
        ``더 올릴 힘''이 약해지고,
      \item ``나쁜 행동''을 샘플링해 자신감을 \textbf{교정}해줄
        기회도 줄어듦.
    \end{itemize}
    $\to$ 단일 샘플 $G_t$ 로 업데이트하는 \textbf{모든} REINFORCE
    계열 알고리즘의 \textbf{본질적 특성}.
  \end{block}
  \vskip0.4em
  {\footnotesize 10.3절에서 가치함수로 추정하는 것(또는 더 많은
  샘플을 쓰는 것)이 이 ``느림''을 공격하는 방법 중 하나.}
  \vskip0.6em
  \begin{block}{왜 이게 model-free인가 --- 다시 한 번}
    궤적 확률 $P(\tau;\theta) = \prod_t \pi_\theta(a_t|s_t)\,
    P(s_{t+1}|s_t,a_t)$ 를 \textbf{로그}로 바꾸면, 환경 전이확률 항은
    $\theta$ 와 \textbf{무관}해서 미분하면 \textbf{사라짐}.
    최종 식에는 정책 $\pi_\theta$ \textbf{만} 남고 환경 모델은
    전혀 등장하지 않음 --- model-free의 핵심 성질이 여기서 나옴.
  \end{block}
\end{frame}
